\chapter{C++ for Latency I: Memory}\label{ch:ll:cpp-for-latency-i-memory}

Writing one byte to a page of memory that has just been mapped cost this book's laptop about \qty{1.3}{\micro\second};
writing the same byte to the same page a second time cost \qty{30}{\nano\second}. The first write went to the kernel,
which found a free physical page, zeroed it and installed it in the page tables; the second was an ordinary store. A
trading process that asks the heap for memory while it handles a message may, now and then, take that path, or several
others inside the allocator that are fast on average and slow once in a while. This chapter is about memory in C++ on
the hot path: how objects are laid out, what an allocation costs and when, and the arenas, pools and \omterm{def:ll:cpp-for-latency-i-memory:sbo}{fixed-capacity
containers} that let the hot path run without allocating at all.

\section{Object layout, padding and alignment}

\begin{definition}[Structure padding]\label{def:ll:cpp-for-latency-i-memory:padding}
Every scalar type has an alignment, equal to its size on x86-64 (a \texttt{double} starts at a multiple of 8). The
compiler lays out a structure's fields in declaration order and inserts \emph{structure padding}\index{structure
padding}, unused bytes, before any field that would otherwise be misaligned and at the end, so that the structure's size
is a multiple of its largest alignment.
\end{definition}

\begin{example}[Ordering fields]\label{ex:ll:cpp-for-latency-i-memory:layout}
An order written as a first draft declares a \texttt{char} side, a \texttt{double} price, a \texttt{char} flag, a 64-bit
id, a 32-bit quantity and a \texttt{char} time in force: 23 bytes of data in 40 bytes, 17 of them padding. Ordered by size
(id, price, quantity, then the three characters) the same fields take 24 bytes: eight orders occupy three 64-byte \omterm{def:ll:memory-hierarchy-and-caches:line}{cache lines}
instead of five. The compiler will not reorder fields itself: the order is part of the type's interface.
\end{example}

\begin{omfigure}
\begin{tikzpicture}[font=\footnotesize,x=0.17cm,y=1cm]
\newcommand{\fld}[4]{\draw[fill=#4] (#1,#3) rectangle (#1+#2,#3+0.55);}
\node[anchor=east] at (-1,0.28) {first draft (40 bytes)};
\fld{0}{1}{0}{omDef!45}\fld{1}{7}{0}{gray!15}\fld{8}{8}{0}{omDef!45}\fld{16}{1}{0}{omDef!45}\fld{17}{7}{0}{gray!15}
\fld{24}{8}{0}{omDef!45}\fld{32}{4}{0}{omDef!45}\fld{36}{1}{0}{omDef!45}\fld{37}{3}{0}{gray!15}
\node at (4,0.28) {pad}; \node at (12,0.28) {price}; \node at (20.5,0.28) {pad}; \node at (28,0.28) {id}; \node at (34,0.28) {qty};
\node[anchor=east] at (-1,-0.72) {by size (24 bytes)};
\fld{0}{8}{-1}{omMeth!40}\fld{8}{8}{-1}{omMeth!40}\fld{16}{4}{-1}{omMeth!40}\fld{20}{3}{-1}{omMeth!40}\fld{23}{1}{-1}{gray!15}
\node at (4,-0.72) {id}; \node at (12,-0.72) {price}; \node at (18,-0.72) {qty};
\foreach \x in {0,8,16,24,32,40,48,56,64} {\draw[gray] (\x,-1.25) -- (\x,-1.1); \node[below,font=\scriptsize] at (\x,-1.2) {\x};}
\draw[omThm,dashed] (64,-1.1) -- (64,0.7);
\node[omThm,anchor=south,font=\scriptsize] at (64,0.7) {end of a 64-byte line};
\end{tikzpicture}

\omcaption{The order of \cref{ex:ll:cpp-for-latency-i-memory:layout} laid out byte by byte: the first draft carries 17
bytes of padding (grey) in 40; the same fields sorted by size need 24, with the side, flag and time in force in the three
bytes after the quantity.}\label{fig:ll:cpp-for-latency-i-memory:layout}
\end{omfigure}

Three further rules. Put the fields the hot path reads together at the start of the structure, so that one line serves the
common case, and move the rest (text, audit fields) to a separate structure. Align structures written by different threads
to separate \omterm{def:ll:memory-hierarchy-and-caches:line}{cache lines} (chapter~3). And prefer indices into arrays to pointers in hot structures: a 32-bit index is half a
pointer and survives a relocation of the array.

\section{What an allocation costs, and when}

\texttt{new} and \texttt{malloc} are fast on average. The GNU C library's allocator keeps small freed blocks in per-thread
caches and hands them back in a few tens of nanoseconds. It also, occasionally, does much more: it consolidates free
blocks, grows the heap with a system call, maps large blocks directly from the kernel (by default at 128~KiB and above, a
threshold it adjusts as the program runs), returns memory to the kernel when the top of the heap is free (by default beyond
128~KiB), and every new page it touches costs a \omterm{def:ll:cpp-for-latency-i-memory:fault}{page fault}.

\begin{definition}[Page fault, pre-faulting]\label{def:ll:cpp-for-latency-i-memory:fault}
A \emph{page fault}\index{page fault} is the processor's trap into the kernel when a program touches a virtual page that
is not yet backed by physical memory (a minor fault: the kernel allocates and zeroes a page) or whose contents are on disk
(a major fault). \emph{Pre-faulting}\index{pre-faulting} is touching every page of a buffer at start-up, so that no fault
happens later on the hot path.
\end{definition}

\Cref{fig:ll:cpp-for-latency-i-memory:faults} measures both costs. Replacing one of 10\,000 live 64-byte orders at random
costs a new-and-delete pair about \qty{25}{\nano\second} at the median, a pool slightly less, and both reach tens of
microseconds at their maximum, where the thread was interrupted. The first write to each page of a fresh 64~MiB mapping
costs over a microsecond a page: 16\,384 faults, forty times the cost of the same writes once the pages exist. A hot path
that allocates is exposed to the allocator's slow paths and to faults on memory it has not touched; one that does not
allocate is exposed to neither.

\begin{omfigure}
\begin{tikzpicture}
\begin{axis}[name=churn,width=7.2cm,height=5.4cm,ybar,bar width=5pt,ymode=log,log origin y=infty,ymin=5,ymax=200000,
  symbolic x coords={new/delete,pool,pmr pool},xtick=data,enlarge x limits=0.3,title={\small replacing one order},
  ylabel={ns (log scale)},tick label style={font=\footnotesize},label style={font=\small},grid=major,grid style={gray!25},
  legend columns=4,legend style={font=\footnotesize,at={(1.15,-0.22)},anchor=north,draw=none,fill=none,
  /tikz/every even column/.append style={column sep=6pt}},table/col sep=comma]
\addplot[fill=omDef!70,draw=omDef] table[x=allocator,y=p50]{figdata/low-latency/06-cpp-for-latency-i-memory/measured_churn.csv};
\addplot[fill=omDef!35,draw=omDef] table[x=allocator,y=p99]{figdata/low-latency/06-cpp-for-latency-i-memory/measured_churn.csv};
\addplot[fill=omThm!45,draw=omThm] table[x=allocator,y=p999]{figdata/low-latency/06-cpp-for-latency-i-memory/measured_churn.csv};
\addplot[fill=gray!40,draw=gray] table[x=allocator,y=max]{figdata/low-latency/06-cpp-for-latency-i-memory/measured_churn.csv};
\legend{p50,p99,p99.9,maximum}
\end{axis}
\begin{axis}[at={(churn.east)},anchor=west,xshift=2.1cm,width=5.2cm,height=5.4cm,ybar,bar width=14pt,ymode=log,log origin y=infty,
  ymin=5,ymax=5000,symbolic x coords={first touch,pre-faulted},xtick=data,enlarge x limits=0.5,title={\small writing a page},
  ylabel={ns per page (log scale)},tick label style={font=\footnotesize},label style={font=\small},grid=major,
  grid style={gray!25},table/col sep=comma]
\addplot[fill=omThm!45,draw=omThm] table[x=kind,y=ns_per_page]{figdata/low-latency/06-cpp-for-latency-i-memory/measured_faults.csv};
\end{axis}
\end{tikzpicture}

\omcaption{Left: destroying and creating one of 10\,000 live 64-byte orders at random with the C++ heap, the build's pool
and the standard library's pool resource (400\,000 replacements). Right: the first write to each 4~KiB page of a fresh
64~MiB mapping, and to the same pages once written. Measured on a laptop (Intel Core Ultra 7 155H) under WSL2, no
isolated cores, pinned to one CPU. Data: \texttt{bench\_alloc.py}.}\label{fig:ll:cpp-for-latency-i-memory:faults}
\end{omfigure}

\section{Arenas, pools and free lists}

\begin{definition}[Arena allocator]\label{def:ll:cpp-for-latency-i-memory:arena}
An \emph{arena allocator}\index{arena allocator} hands out consecutive pieces of one pre-allocated block by advancing an
offset, and frees them all at once by resetting it; individual objects are never freed.
\end{definition}

\begin{definition}[Object pool, free list]\label{def:ll:cpp-for-latency-i-memory:pool}
An \emph{object pool}\index{object pool} pre-allocates storage for a fixed number of objects of one type and recycles it:
creating takes a free slot, destroying returns it. A \emph{free list}\index{free list} is the list of free slots; an
intrusive free list stores its links inside the free slots themselves, so it needs no memory of its own.
\end{definition}

An arena suits memory whose lifetime is a phase: the scratch space of one message, rebuilt for the next; the parsed fields
of a batch. A pool suits objects of one type with independent lifetimes: orders, which are created and cancelled in any
order. Both have fixed capacity decided in advance, and the design question they force is the right one: what happens when
it runs out. On the hot path the answer is never ``ask the heap''; it is to refuse (reject the order, drop to a degraded
mode) and raise an alarm, and to size the pool from the measured peak with a margin.

The standard library has the same ideas as memory resources (\texttt{std::pmr}, since C++17): a monotonic buffer resource
is an arena, a pool resource keeps \omterm{def:ll:cpp-for-latency-i-memory:pool}{free lists} by size. They let standard containers use them, at the price of an indirect
call per allocation; the general-purpose pool of \cref{fig:ll:cpp-for-latency-i-memory:faults} is slower than a pool for one
type.

\section{Small buffers and fixed-capacity containers}

\begin{definition}[Small-buffer optimisation, fixed-capacity container]\label{def:ll:cpp-for-latency-i-memory:sbo}
A type with a \emph{small-buffer optimisation}\index{small-buffer optimisation} stores small contents inside the object and
allocates only for large ones: the GNU implementation of \texttt{std::string} holds up to 15 characters inline. A
\emph{fixed-capacity container}\index{fixed-capacity container} always stores its elements inside the object, up to a
capacity fixed at compile time, and refuses to grow beyond it.
\end{definition}

The \omterm{def:ll:cpp-for-latency-i-memory:sbo}{small-buffer optimisation} makes allocation depend on the data: a symbol of four letters does not allocate, a client
order identifier of 22 characters does, and a test with short identifiers will not see it. The handler of
\cref{lst:ll:cpp-for-latency-i-memory:naive} is typical of first drafts. Counted with a replaced global \texttt{operator
new}, it allocates nine times per message: one for the long identifier, three as the vector of fills grows to three, two
for the new map entry (the tree node and its copy of the identifier), one when the lambda copies the fills, and two when
\texttt{std::function} stores the lambda on the heap and copies it. The rewrite with a fixed string, a fixed vector of
fills and a fixed table of entries does the same work with none.

\begin{method}[Removing allocation from a hot path]\label{met:ll:cpp-for-latency-i-memory:remove}
\begin{enumerate}
\item Count: replace the global \texttt{operator new} in a test build and assert the count per message after warm-up.
\item For each allocation, decide the capacity: identifiers and symbols in fixed strings, per-message lists in fixed
vectors, objects with lifetimes in pools, per-message scratch in an arena reset per message.
\item Replace type-erased callables that allocate (\texttt{std::function} with large captures) by templates (chapter~7) or
function pointers with a context.
\item Allocate and touch every pool, arena and table at start-up, on the thread that will use it (chapter~4).
\item Keep the assertion in the tests, so that a later change that allocates fails the build.
\end{enumerate}
\end{method}

\section{Tutorial: counting and removing allocations}\label{tut:ll:cpp-for-latency-i-memory:tut}

\begin{tutorial}
\textbf{Goal.} Count the allocations of a message handler, remove them, and measure what allocation and \omterm{def:ll:cpp-for-latency-i-memory:fault}{page faults} cost.
\textbf{End state:} a handler with zero allocations per message after warm-up, and
\cref{fig:ll:cpp-for-latency-i-memory:faults}.
\begin{tutsteps}
\item \textbf{The first draft.} Strings, a vector that grows, a map, a callback.
\omcode{code/low-latency/06-cpp-for-latency-i-memory/cpp/ll_alloc.hpp}{22}{35}{A message handler that allocates nine times per message.}\label{lst:ll:cpp-for-latency-i-memory:naive}
\item \textbf{Count.} \texttt{firm\_alloc\_count.hpp} replaces the global \texttt{operator new} with one that increments an
atomic counter; the test handles a message, reads the counter, and asserts 9 for the draft and 0 for the rewrite.
\item \textbf{A pool with an intrusive \omterm{def:ll:cpp-for-latency-i-memory:pool}{free list}.} Each free slot holds the pointer to the next free slot; creating pops
the list, destroying pushes.
\omcode{code/firm/arena/cpp/firm_arena.hpp}{45}{80}{The object pool.}\label{lst:ll:cpp-for-latency-i-memory:pool}
\item \textbf{The same guarantee in Rust.} A counting global allocator (\texttt{firm\_arena::CountingAlloc}) installed
in the test binary asserts that the pool's operations allocate nothing.
\item \textbf{Measure} with \texttt{python bench\_alloc.py}: allocator churn and \omterm{def:ll:cpp-for-latency-i-memory:fault}{page faults}.
\end{tutsteps}
\textbf{What to change next.} Shorten the client identifier to 15 characters and count again; make the pool's slots
cache-line aligned and compare the churn percentiles.
\end{tutorial}

\section{Build: the allocation-free toolkit}\label{bld:ll:cpp-for-latency-i-memory:arena}

\begin{build}
\bfield{Purpose} The storage every hot-path component of Part~IV uses (the order pool of the book builder, the gateway's
order states, per-message scratch), and the allocation counter every acceptance test asserts with.
\bfield{Interface} C++20 \texttt{firm::arena}: \texttt{Arena(bytes)} with \texttt{allocate(n, align)}, \texttt{make<T>},
\texttt{reset}; \texttt{Pool<T, N>} with \texttt{create}, \texttt{destroy}, \texttt{live}; \texttt{FixedVector<T, N>};
\texttt{FixedString<N>}; \texttt{firm\_alloc\_count.hpp} with \texttt{allocations()}. Rust \texttt{firm\_arena}:
\texttt{Arena} (offsets), \texttt{Pool<T>} (indices), \texttt{CountingAlloc} and \texttt{allocations()}.
\bfield{Rules} Capacities are fixed at construction; exhaustion returns a null result, never a heap allocation; storage is
pre-faulted at construction; no operation after construction allocates.
\bfield{Acceptance tests} \texttt{code/firm/arena/}: aligned bumps, refusal when full and reset; the pool's reuse order and
exhaustion; fixed strings; zero allocations counted around all of it, in C++20 and Rust.
\bfield{Stretch} A pool whose slots are aligned to \omterm{def:ll:memory-hierarchy-and-caches:line}{cache lines} for objects written by different threads; an arena that
records its high-water mark for sizing.
\end{build}

\begin{omsources}
\item Linux man page \texttt{mallopt(3)}: \texttt{M\_MMAP\_THRESHOLD}, \texttt{M\_TRIM\_THRESHOLD}.
\item cppreference, ``\texttt{std::pmr::memory\_resource}'' and its derived resources.
\item C.~Cook, ``When a microsecond is an eternity'', CppCon 2017 (reusing memory on the hot path).
\end{omsources}

\section{Exercises}

\begin{exercise}[$\star$]\label{exo:ll:cpp-for-latency-i-memory:1}
What are the size and padding of a structure holding, in this order, a \texttt{char}, an \texttt{int32}, a \texttt{char},
a \texttt{double} and an \texttt{int16}? And with the fields sorted by size?
\end{exercise}

\begin{exercise}[$\star$]\label{exo:ll:cpp-for-latency-i-memory:2}
A pool of 50\,000 orders of 64 bytes: how much memory, and how many 4~KiB pages must be pre-faulted?
\end{exercise}

\begin{exercise}[$\star$]\label{exo:ll:cpp-for-latency-i-memory:3}
At the measured cost per page, how long does touching a fresh 512~MiB buffer take at start-up, and what would it cost if
it happened on the hot path instead, a page at a time?
\end{exercise}

\begin{exercise}[$\star\star$]\label{exo:ll:cpp-for-latency-i-memory:4}
Which of these allocate with the GNU standard library: \texttt{std::string s(\textquotedbl{}AAPL\textquotedbl{})};
\texttt{std::string s(\textquotedbl{}CLIENT-ORDER-ID-42\textquotedbl{})}; \texttt{std::vector<int> v; v.reserve(0)};
\texttt{std::vector<int> v(1)}?
\end{exercise}

\begin{exercise}[$\star\star$]\label{exo:ll:cpp-for-latency-i-memory:5}
The draft handler allocates nine times per message. Which three changes remove the most allocations, and how many does each
remove?
\end{exercise}

\begin{exercise}[$\star\star$]\label{exo:ll:cpp-for-latency-i-memory:6}
Read \cref{fig:ll:cpp-for-latency-i-memory:faults}: at the median the heap is almost as fast as the pool. Why use a pool
at all?
\end{exercise}

\begin{exercise}[$\star\star\star$]\label{exo:ll:cpp-for-latency-i-memory:7}
\emph{Coding.} Give \texttt{firm::arena::Pool} a capacity check that records the high-water mark of live objects, run the
churn benchmark with it, and propose a sizing rule for a pool whose peak is measured over a year of busy days.
\end{exercise}

\begin{exercise}[$\star\star\star$]\label{exo:ll:cpp-for-latency-i-memory:8}
\emph{Find the flaw.} ``Our handler allocates nothing: we checked it with a test that sends the same message, symbol
\texttt{ESZ6} and client id \texttt{A1}, a million times.''
\end{exercise}

\section{Problem: Four Hundred Microseconds, Once a Minute}

\begin{problem}[{Weekend problem --- a handler that allocated}]\label{pb:ll:cpp-for-latency-i-memory:1}
A strategy's handler takes about \qty{2}{\micro\second} a message, except once or twice a minute when it takes several
hundred. The team suspects memory. The handler is the draft of \cref{lst:ll:cpp-for-latency-i-memory:naive}, fed 50\,000
messages a second.

\medskip\noindent\textbf{Part I --- Counting.}
\begin{enumerate}
\item How many allocations a second does the draft make?
\item List the nine allocations of one message and the line that makes each.
\item Which of them would disappear with an identifier of 15 characters or fewer?
\item Why does the count not depend on how many fills a message has, beyond three?
\end{enumerate}

\medskip\noindent\textbf{Part II --- The costs.}
\begin{enumerate}[resume]
\item From \cref{fig:ll:cpp-for-latency-i-memory:faults}, what does a pair of heap operations cost at the median?
\item How much time a second do the draft's allocations cost at that median?
\item What does a \omterm{def:ll:cpp-for-latency-i-memory:fault}{page fault} cost, and how many would put a message at several hundred microseconds?
\item Name three events inside the allocator that are rare but slow.
\end{enumerate}

\medskip\noindent\textbf{Part III --- The rewrite.}
\begin{enumerate}[resume]
\item What replaces each allocation in the fixed handler?
\item What capacities did the rewrite choose, and what happens when one is exceeded?
\item Why must the fixed handler's storage be touched at start-up, and by which thread?
\item What does the rewrite's linear search over 64 entries cost, and when would it need a hash table instead?
\end{enumerate}

\medskip\noindent\textbf{Part IV --- The verdict.}
\begin{enumerate}[resume]
\item State the \emph{named result}: allocations per message before and after the rewrite, and the cost of a \omterm{def:ll:cpp-for-latency-i-memory:fault}{page fault}
against a pre-faulted write on this laptop.
\item How does the test keep the result from regressing?
\item What else, besides allocation, can make one message in a minute slow?
\item How would you show which cause it was on the production host?
\item Why is ``refuse when full'' acceptable on the hot path, and ``allocate when full'' not?
\item Where does the Rust twin enforce the same rule?
\item What should the pool's size be if the peak of live orders in a year was 38\,000?
\item In one sentence: what is the memory rule of a hot path?
\end{enumerate}
\end{problem}

\section{Interview questions}

\begin{interviewq}[$\star$ \iqroles{developer}]\label{iq:ll:cpp-for-latency-i-memory:1}
Why can reordering the fields of a structure make a program faster?
\end{interviewq}

\begin{interviewq}[$\star\star$ \iqroles{developer}]\label{iq:ll:cpp-for-latency-i-memory:2}
Why do low-latency systems avoid heap allocation on the hot path, if \texttt{malloc} takes 20 nanoseconds?
\end{interviewq}

\begin{interviewq}[$\star\star$ \iqroles{developer}]\label{iq:ll:cpp-for-latency-i-memory:3}
Implement an \omterm{def:ll:cpp-for-latency-i-memory:pool}{object pool} for a fixed-size type. What happens when it is empty?
\end{interviewq}

\begin{interviewq}[$\star\star$ \iqroles{developer}]\label{iq:ll:cpp-for-latency-i-memory:4}
What is the small-string optimisation, and why can it hide allocations from a test?
\end{interviewq}

\begin{interviewq}[$\star\star$ \iqroles{developer}]\label{iq:ll:cpp-for-latency-i-memory:5}
What is a \omterm{def:ll:cpp-for-latency-i-memory:fault}{page fault}, and how do you keep them off a hot path?
\end{interviewq}

\begin{interviewq}[$\star\star\star$ \iqroles{developer}]\label{iq:ll:cpp-for-latency-i-memory:6}
How would you prove, in continuous integration, that a component never allocates after start-up?
\end{interviewq}
